\documentclass[11pt,a4paper]{article}

\setlength{\hoffset}{1in}
\setlength{\oddsidemargin}{-1in}
\setlength{\voffset}{0.9in}
\setlength{\topmargin}{-1.6in}
\setlength{\textwidth}{15.92cm}
\setlength{\textheight}{25cm}
\linespread{1.1}

\usepackage[utf8]{inputenc}
\usepackage{amsmath}
\usepackage{amsfonts}
\usepackage{amssymb}
\usepackage{makeidx}


\begin{document}

\begin{center}
SHENZHEN UNIVERSITY \\
INSTITUTE FOR ADVANCED STUDY \\
 Differential Equations, Autumn 2026 \\
Assignment 3, due  on Thursday, November 9, 2026
\end{center}


\noindent
{\bf Problem 1.} (Problem 7 of Exercises 2.6)
Solve the ODE $\cos (x+y)dy=\sin (x+y)dx$.
\par\medskip
\noindent\textbf{Solution.}\par
\medskip

Let \(u=x+y\).
\[
du=dx+dy
\]
\[
\frac{du}{dx}=1+\frac{dy}{dx}
\]

\[
\cos(x+y)\,dy=\sin(x+y)\,dx
\]
\[
\frac{\cos u}{\sin u}\cdot
\left(\frac{du}{dx}-1\right)=1
\]
\[
\frac{du}{dx}=\frac{\sin u}{\cos u}+1
\]
\[
\frac{du}{dx}=\frac{\sin u+\cos u}{\cos u}
\]
\[
dx=\frac{\cos u}{\sin u+\cos u}\,du
\]

Suppose \(\cos u=A(\sin u+\cos u)+B(\cos u-\sin u)\).
\[
\cos u=(A-B)\sin u+(A+B)\cos u
\]
\[
\begin{cases}
A-B=0,\\
A+B=1
\end{cases}
\quad\Longrightarrow\quad
\begin{cases}
A=\dfrac12,\\[4pt]
B=\dfrac12
\end{cases}
\]
\[
dx=\frac12\left(
\frac{\sin u+\cos u}{\sin u+\cos u}
+\frac{\cos u-\sin u}{\sin u+\cos u}
\right)\,du
\]
\[
\int dx=\frac12\int\left(
1+\frac{\cos u-\sin u}{\sin u+\cos u}
\right)\,du
\]
\[
x=\frac12\left(u+\ln|\sin u+\cos u|\right)+C
\]

Substitute \(u=x+y\).
\[
x=\frac12\left(
x+y+\ln|\sin(x+y)+\cos(x+y)|
\right)+C
\]
\[
x-y-\ln|\sin(x+y)+\cos(x+y)|=C_2,
\qquad C_2\in\mathbb R
\]
\[
\sin u+\cos u=0
\quad\Longrightarrow\quad
y=-x-\frac{\pi}{4}+k\pi,
\qquad k\in\mathbb Z
\]

\bigskip
\bigskip
\noindent
{\bf Problem 2.} (Problem 47 of Exercises 2.6) {{\bf Riccati Equation}}. An equation of the form
\[
\dfrac{dy}{dx}=P(x)y^2+Q(x)y+R(x)
\]
is called a  {\bf Riccati equation }.

a) If one solution $u(x)$ of the above equation is known, show that the substitution $y=u+1/v$ reduces the above equation to a linear equation in $v$ .
\par\medskip
\noindent\textbf{Solution.}\par
\medskip
Substitute \(y=u+\dfrac1v\).
\[
\frac{du}{dx}-\frac1{v^2}\frac{dv}{dx}
=P(x)\left(u+\frac1v\right)^2
+Q(x)\left(u+\frac1v\right)+R(x)
\]
\[
\frac{du}{dx}-\frac1{v^2}\frac{dv}{dx}
=P(x)\left(u^2+\frac{2u}{v}+\frac1{v^2}\right)
+Q(x)\left(u+\frac1v\right)+R(x)
\]
\[
\begin{aligned}
\frac{du}{dx}-\frac1{v^2}\frac{dv}{dx}
&=\left(P(x)u^2+Q(x)u+R(x)\right)\\
&\quad+\frac1v\left(2P(x)u+Q(x)\right)
+\frac{P(x)}{v^2}
\end{aligned}
\]
\[
-\frac1{v^2}\frac{dv}{dx}
=\frac1v\left(2P(x)u+Q(x)\right)
+\frac{P(x)}{v^2}
\]
\[
\frac{dv}{dx}
=-v\left(2P(x)u+Q(x)\right)-P(x)
\]
\[
\frac{dv}{dx}+\left(2P(x)u+Q(x)\right)v=-P(x)
\]

\medskip
b) Given that $u(x)=x$ is a solution to $\dfrac {dy}{dx}=x^3(y-x)^2+\dfrac {y}{x} \quad$,
use the result of part (a) to find all the other solutions to this equation.
\medskip
\par\medskip
\noindent\textbf{Solution.}\par
\medskip
\[
\frac{dy}{dx}=x^3(y-x)^2+\frac yx
\]
\[
\frac{dy}{dx}=x^3(y^2-2xy+x^2)+\frac yx
\]
\[
\frac{dy}{dx}=x^3y^2-2x^4y+x^5+\frac yx
\]
\[
\frac{dy}{dx}
=x^3y^2+\left(\frac1x-2x^4\right)y+x^5
\]
\[
P(x)=x^3,\qquad
Q(x)=\frac1x-2x^4,\qquad
R(x)=x^5
\]

Let \(y=u+\dfrac1v=x+\dfrac1v\).
\[
\frac{dv}{dx}
+\left(2x^3\cdot x+\frac1x-2x^4\right)v=-x^3
\]
\[
\frac{dv}{dx}+\frac vx=-x^3
\]
\[
x\frac{dv}{dx}+v=-x^4
\]
\[
\frac{d}{dx}(xv)=-x^4
\]
\[
\int\frac{d}{dx}(xv)\,dx=-\int x^4\,dx
\]
\[
xv=-\frac{x^5}{5}+C
\]
\[
v=-\frac{x^4}{5}+\frac Cx
\]

Substitute \(v\) into \(y=u+\dfrac1v\).
\[
y=x+\frac{5x}{C_1-x^5},
\qquad \text{while }y=x\text{ is also a solution}
\]

\bigskip
\noindent
{\bf Problem 3.} (Problem 33 of Exercises 2.4)
Find the {\bf orthogonal trajectories} to the families (refer to Problem 32 on page 62 and 63 for the definition of  {\bf orthogonal trajectories})

a)  $2x^2 + y^2 = k$; \;\;\;\; b) $y = k x^4$.
\par\medskip
\noindent\textbf{Solution.}\par
\medskip

\medskip
\noindent a)
\[
2x^2+y^2=k
\]
\[
4x\,dx+2y\,dy=0
\]
\[
\frac{dy}{dx}=-\frac{2x}{y}
\]

The orthogonal trajectories should satisfy this ODE:
\[
\frac{dy}{dx}=\frac{y}{2x}
\]
\[
\frac{dy}{y}=\frac{dx}{2x}
\]
\[
\ln|y|=\frac12\ln|2x|+C
\]
\[
y^2=C_1x
\]

\medskip
\noindent b)
\[
y=kx^4
\]
\[
\frac{dy}{dx}=4kx^3=\frac{4y}{x}
\]

The orthogonal trajectories should satisfy this ODE:
\[
\frac{dy}{dx}=-\frac{x}{4y}
\]
\[
4y\,dy=-x\,dx
\]
\[
\int4y\,dy=-\int x\,dx
\]
\[
2y^2=-\frac12x^2+C
\]
\[
x^2+4y^2=C_1
\]

\bigskip
\bigskip
\noindent
{\bf Problem 4.} (Problem 5 of Exercises 3.2) A swimming pool whose volume is 10,000 gal contains water
that is $0.01\%$ chlorine. Starting at $t=0$, city water containing
$0.001\%$ chlorine is pumped into the pool at a rate of 5 gal/min. The pool water flows out at the same rate.
What is the percentage of chlorine in the pool after 1 h? When will the pool water be $0.002\%$ chlorine?
\par\medskip
\noindent\textbf{Solution.}\par
\medskip

Suppose \(y(t)\) is the concentration of chlorine at the time \(t\)
in the swimming pool.
\[
10^4\frac{dy}{dt}=5\times10^{-5}-5y
\]
\[
\frac{dy}{5\times10^{-9}-5\times10^{-4}y}=dt
\]
\[
\frac1{5\times10^{-4}}
\int\frac{dy}{10^{-5}-y}
=\int dt
\]
\[
-2000\ln|10^{-5}-y|=t+C
\]
\[
\ln|10^{-5}-y|=-\frac1{2000}t+C_1
\]
\[
y=10^{-5}+C_2e^{-t/2000}
\]

Substitute \(y(0)=10^{-4}\).
\[
10^{-4}=10^{-5}+C_2
\]
\[
C_2=0.00009
\]
\[
y(t)=10^{-5}+0.00009e^{-t/2000}
\]

\[
y(60)\approx0.0000973401
\]
\[
y(60)\times100\%\approx0.009734\%
\]

\[
y(t)=0.00002
\]
\[
t\approx4394.45\ \mathrm{min}\approx73.24\ \mathrm{h}
\]

\bigskip


\bigskip
\noindent
{\bf Problem 5.} (Problem 1, Exercises 3.3) A cup of hot coffee initially at $95^{\circ}$C cools to $80^{\circ}$C in 5 minutes while
sitting in a room of temperature $21^{\circ}$C. Using just Newton's law of cooling, determine when the temperature of the coffee will be
a nice $50^{\circ}$C.
\par\medskip
\noindent\textbf{Solution.}\par
\medskip

\[
T(t)=T_{\mathrm{out}}+(T_0-T_{\mathrm{out}})e^{-kt}
\]
\[
T(t)=21+(95-21)e^{-kt}
\]
\[
T(t)=21+74e^{-kt}
\]

Substitute \(T(5)=80\).
\[
80=21+74e^{-5k}
\]
\[
74e^{-5k}=59
\]
\[
e^{-5k}=\frac{59}{74}
\]
\[
-5k=\ln\frac{59}{74}
\]
\[
k=-\frac15\ln\frac{59}{74}
\]
\[
T(t)=21+74e^{\frac15\ln\frac{59}{74}\,t}
\]

\[
T(t)=50
\]
\[
t\approx20.68\ \mathrm{min}
\]

\bigskip

\bigskip
\noindent
{\bf Problem 6.} (Problem 7, Exercises 3.3) In a hot Saturday morning while people are working inside, the air conditioner
keeps the temperature inside the building at $24^{\circ}$C. At noon the air conditioner is turned off, and the people
go home. The temperature outside is a constant $35^{\circ}$C for the rest of the afternoon. If the time constant for the building is 4hr, what will
be the temperature inside the building at 2:00pm? At 6:00pm? When will the
temperature inside the building reach $27^{\circ}$C?
\par\medskip
\noindent\textbf{Solution.}\par
\medskip

\[
T(t)=T_{\mathrm{out}}
+(T_0-T_{\mathrm{out}})e^{-t/\tau}
\]
\[
T(t)=35-11e^{-t/4}
\]
\[
T(2)=35-11e^{-1/2}\approx28.3^\circ\mathrm C
\]
\[
T(6)=35-11e^{-3/2}\approx32.5^\circ\mathrm C
\]
\[
T(t)=27
\]
\[
t\approx1.2738\ \mathrm h,
\qquad \text{time: 1:16 p.m.}
\]

\bigskip

\bigskip
\noindent
{\bf Problem 7.} (Problem 5 of Exercises 3.4) An object  of mass 5 kg is given an initial downward velocity of 50m/sec and then allowed to fall
under the influence of gravity. Assume that the force in newtons due to air resistance is $-10v$, where $v$ is
the velocity of the object in m/sec. Determine the equation of motion of the object. If the object is initially
500 m above the ground, determine when the object will strike the ground.
\par\medskip
\noindent\textbf{Solution.}\par
\medskip

\[
m\frac{dv}{dt}=mg-10v
\]
\[
5\frac{dv}{dt}=5\times9.8-10v
\]
\[
\frac{dv}{49-10v}=\frac15\,dt
\]
\[
\int\frac{dv}{49-10v}=\frac15\int dt
\]
\[
-\frac1{10}\ln|49-10v|=\frac15t+C
\]
\[
\ln|49-10v|=-2t+C_1
\]
\[
49-10v=C_2e^{-2t}
\]

Substitute \(v(0)=50\).
\[
49-500=C_2
\]
\[
C_2=-451
\]

\[
49-10v=-451e^{-2t}
\]
\[
v=\frac{49}{10}+\frac{451}{10}e^{-2t}
\]
\[
x=500-\int v\,dt
=500-\left(
\frac{49}{10}t-\frac{451}{20}e^{-2t}+C_3
\right)
\]

Substitute \(x(0)=500\).
\[
500=500-\left(-\frac{451}{20}+C_3\right)
\]
\[
C_3=\frac{451}{20}
\]
\[
x(t)=500-\frac{49}{10}t
+\frac{451}{20}e^{-2t}-\frac{451}{20}
\]
\[
x(t)=\frac{9549}{20}-\frac{49}{10}t
+\frac{451}{20}e^{-2t}
\]

\[
x(t)=0
\]
\[
t\approx97.44\ \mathrm s
\]

\bigskip
\bigskip
\noindent
{\bf Problem 8.}
Find the solutions to the following ODEs.

a) $y''+6y'+9y = 0$; \;\;\;\;\; b) $6y''+y'-2y = 0$.
\par\medskip
\noindent\textbf{Solution.}\par
\medskip

\medskip
\noindent a)
\[
y''+6y'+9y=0
\]

Suppose \(y=e^{\alpha t}\).
\[
y''+6y'+9y=(\alpha^2+6\alpha+9)Ae^{\alpha t}=0
\]
\[
\alpha^2+6\alpha+9=0
\]
\[
(\alpha+3)^2=0
\]
\[
\alpha=-3
\]
\[
y_1=e^{-3t},\qquad y_2=te^{-3t}
\]
\[
y=C_1y_1+C_2y_2=(C_1+C_2t)e^{-3t},
\qquad C_1,C_2\in\mathbb R
\]

\medskip
\noindent b)
\[
6y''+y'-2y=0
\]

Suppose \(y=e^{\alpha t}\).
\[
6\alpha^2+\alpha-2=0
\]
\[
(3\alpha+2)(2\alpha-1)=0
\]
\[
\alpha_1=-\frac23,\qquad \alpha_2=\frac12
\]
\[
y=C_1e^{-\frac23t}+C_2e^{t/2},
\qquad C_1,C_2\in\mathbb R
\]

\bigskip
\noindent
{\bf Problem 9.}
Solve the following initial value problems.

a) $y''+2y'-8y = 0 ; y(0)=3, \;\;\;y'(0)= -12$;

b) $y''-4y'-5y = 0 ; y(-1)= 3,\;\;\; y'(-1) = 9$.

\par\medskip
\noindent\textbf{Solution.}\par
\medskip

\medskip
\noindent a)
\[
y''+2y'-8y=0
\]

Suppose \(y=e^{\alpha t}\).
\[
\alpha^2+2\alpha-8=0
\]
\[
(\alpha+4)(\alpha-2)=0
\]
\[
\alpha_1=-4,\qquad \alpha_2=2
\]
\[
y=C_1e^{-4t}+C_2e^{2t},
\qquad
y'=-4C_1e^{-4t}+2C_2e^{2t}
\]

Substitute \(y(0)=3\) and \(y'(0)=-12\).
\[
\begin{pmatrix}
1 & 1\\
-4 & 2
\end{pmatrix}
\begin{pmatrix}
C_1\\
C_2
\end{pmatrix}
=
\begin{pmatrix}
3\\
-12
\end{pmatrix}
\]
\[
\begin{pmatrix}
C_1\\
C_2
\end{pmatrix}
=
\begin{pmatrix}
3\\
0
\end{pmatrix}
\]
\[
y=3e^{-4t}
\]

\medskip
\noindent b)
\[
y''-4y'-5y=0
\]

Suppose \(y=e^{\alpha t}\).
\[
\alpha^2-4\alpha-5=0
\]
\[
(\alpha-5)(\alpha+1)=0
\]
\[
\alpha_1=5,\qquad \alpha_2=-1
\]
\[
y=C_1e^{5t}+C_2e^{-t},
\qquad
y'=5C_1e^{5t}-C_2e^{-t}
\]

Substitute \(y(-1)=3\), \(y'(-1)=9\).
\[
\begin{pmatrix}
e^{-5} & e\\
5e^{-5} & -e
\end{pmatrix}
\begin{pmatrix}
C_1\\
C_2
\end{pmatrix}
=
\begin{pmatrix}
3\\
9
\end{pmatrix}
\]
\[
\begin{pmatrix}
C_1\\
C_2
\end{pmatrix}
=
\begin{pmatrix}
2e^5\\
e^{-1}
\end{pmatrix}
\]
\[
y=2e^{5+5t}+e^{-t-1}
\]

\bigskip

\noindent {\bf Problem 10.} (Problems 10 and 20 of Exercises 4.4) Find a particular solution to the given
differential equation

1) $y'' + 2y' - y = 10$;
\medskip
\par\medskip
\noindent\textbf{Solution.}\par
\medskip
\[
y''+2y'-y=10
\]

Notice that \(y=-10\) is a particular solution.

\medskip
2) $y'' + 4y = 16t \sin 2t$.
\par\medskip
\noindent\textbf{Solution.}\par
\medskip
\[
y''+4y=16t\sin2t
\]

Suppose \(y=(At^2+Bt)\sin2t+(Ct^2+Dt)\cos2t\).
\[
\begin{aligned}
y'
&=(2At+B)\sin2t+2\cos2t(At^2+Bt)
 +(2Ct+D)\cos2t\\
&\quad-2\sin2t(Ct^2+Dt)\\
&=\left[-2Ct^2+2(A-D)t+B\right]\sin2t\\
&\quad+\left[2At^2+2(B+C)t+D\right]\cos2t
\end{aligned}
\]
\[
\begin{aligned}
y''
&=\left[-4At^2-(4B+8C)t+2A-4D\right]\sin2t\\
&\quad+\left[-4Ct^2+(8A-4D)t+4B+2C\right]\cos2t
\end{aligned}
\]
\[
y''+4y
=(-8Ct+2A-4D)\sin2t+(8At+4B+2C)\cos2t
\]
\[
\begin{cases}
-8C=16,\\
2A-4D=0,\\
8A=0,\\
4B+2C=0
\end{cases}
\quad\Longrightarrow\quad
\begin{cases}
A=0,\\
B=1,\\
C=-2,\\
D=0
\end{cases}
\]
\[
y=t\sin2t-2t^2\cos2t
\]

\bigskip
\end{document}












